% -----------------------------------------------------------
\section{Goodness, Good}
% -----------------------------------------------------------
	\label{GOODNESS}
	
% % ----------------------
% \subsection{See also}
% % ----------------------

% % ----------------------
% \subsection{1828 American Dictionary of the English Language} % *1828
% % ----------------------
	% \label{GOODNESS-1828}

	% \begin{compactitem}
		% \item
	% \end{compactitem}

% % ----------------------
% \subsection{Oxford English Dictionary} % *OED
% % ----------------------
	% \label{GOODNESS-OED}

	% \begin{compactitem}
		% \item[]
	% \end{compactitem}

% ----------------------
\subsection{Scriptural Usage} % *SCRIPTURES
% ----------------------
	\label{GOODNESS-SCRIPTURES}

	% % -----
	% \subsubsection{Old Testament} % *OT
	% % -----
		% \label{GOODNESS-OT}
	
		% % --
		% \paragraph{KJV}
		% % --
			% \label{GOODNESS-OT-KJV}
	
			% \begin{tabular}{ll}
				% Total occurrences in the OT & 
			% \end{tabular}
			
			% \begin{compactdesc}
				% \item[]
			% \end{compactdesc}
			
		% % --
		% \paragraph{Hebrew Bible}
		% % --
			% \label{GOODNESS-HEBREW}
		
			% %
			% \subparagraph{\texorpdfstring{\he{sample word}}{}}
			% %
				% \label{PAGA} % whatever is in step bible without periods
			
				% \begin{compactdesc}
					% \item[HALOT , PDF ] \textbf{}
						% \begin{compactitem}
							% \item[1]
						% \end{compactitem}
					% \item[BDB , PDF ] \textbf{}
						% \begin{compactitem}
							% \item[1]
						% \end{compactitem}
					% \item[DCH PDF ] \textbf{}
						% \begin{compactitem}
							% \item[1]
						% \end{compactitem}
				% \end{compactdesc}
				
				% \begin{compactdesc}
					% \item[Some OT uses] \textbf{}
						% \begin{compactdesc}
							% \item[]
						% \end{compactdesc}
				% \end{compactdesc}
			
		% % --
		% \paragraph{Septuagint}
		% % --
			% \label{GOODNESS-LXX}
		
			% %
			% \subparagraph{\texorpdfstring{\gr{sample word}}{}}
			% %
		
	% -----
	\subsubsection{New Testament} % *NT
	% -----
		\label{GOODNESS-NT}
	
		% % --
		% \paragraph{KJV}
		% % --
			% \label{GOODNESS-NT-KJV}
			
	
			% \begin{tabular}{ll}
				% Total occurrences in the NT & 
			% \end{tabular}
			
			% \begin{compactdesc}
				% \item[]
			% \end{compactdesc}
			
		% --
		\paragraph{Greek New Testament} % *GREEK
		% --
			\label{GOODNESS-GREEK}
			
			%
			\subparagraph{\texorpdfstring{\gr{ἀγαθός}}{}}
			%
				\label{AGATHOS}
			
				\begin{compactdesc}
					\item[BDAG 2b] \textbf{}
						\begin{compactitem}
							\item[1] \textbf{pertaining to meeting a relatively high standard of quality}, of things
							\item[1.a] adjective \textbf{\textit{useful, beneficial}}
							\item[1.b] used as a pure substantive
							\item[1.b.$\alpha$] quite generally 
								\begin{compactitem}
									\item I think it just means of ``things'' generally; it cites the use in \hyperref[LUKE-16-25]{Luke 16:25}.
								\end{compactitem}
							\item[1.b.$\beta$] \textit{possessions, treasures}
							\item[1.b.$\gamma$] possessions of a higher order
							\item[2] \textbf{pertaining to meeting a high standard of worth and merit, \textit{good}}
							\item[2.a] as adjective
							\item[2.a.$\alpha$] of humans and deities (the primary focus is on usefulness to humans and society in general)
							\item[2.a.$\beta$] of things characterized especially in terms of social significance and worth
							\item[2.b] as substantive, singular
							\item[2.b.$\alpha$] that which is beneficial or helpful
							\item[2.b.$\beta$] \gr{τὰ ἀγαθὰ} \textit{good deeds} (\hyperref[JOHN-5-29]{John 5:29})
						\end{compactitem}
					\item[CGL 3b, PDF 29] \textbf{}
						\begin{compactitem}
							\item[1] (of persons, especially in the performance of a specific role or duty) \textbf{good, capable}
							\item[2] (with connotation of high birth) \textbf{noble}; (with connotation of morality) \textbf{good}
							\item[3] (of persons) \textbf{trustworthy, loyal} (to family, friends, political allies, or similar)
							\item[4] (of thoughts, words, advice, or similar) \textbf{good, useful}
							\item[5] (of a deity) without wickedness, \textbf{beneficent, benevolent}
							\item[6] (as the title of a deity, to whom a toast was made at the end of a meal) \textbf{good}
							\item[7] (ironically, of a ruler) \textbf{esteemed, worthy}
							\item[8] (in vocative address, friendly or sarcastic) \textit{my good man}
						\end{compactitem}
					\item[LSJ 4b, PDF 77] \textbf{}
						\begin{compactitem}
							\item[I] of persons
							\item[I.1] \textit{well-born, gentle}
							\item[I.2] \textit{brave, valiant}, since courage was attributed to Chiefs and Nobles
							\item[I.3] \textit{good, capable}, in reference to ability
							\item[I.4] \textit{good}, in moral sense
							\item[I.5] \textit{my good friend}, as a term of gentle remonstrance
							\item[II] of things
							\item[II.1] \textit{good, serviceable}
							\item[II.2] of outward circumstances
							\item[II.3] \textit{morally good}
								\begin{compactitem}
									\item It lists \hyperref[ROMANS-2-7]{Romans 2:7} as an example.
								\end{compactitem}
							\item[II.4] neuter substantive, \textit{good, blessing, benefit}, of persons or things
							\item[III] comparative and superlative are usually supplied from other stems
							\item[IV] adverb usually \gr{εὖ}
						\end{compactitem}
				\end{compactdesc}
		
			%
			\subparagraph{\texorpdfstring{\gr{καλός}}{}}
			%
				\label{KALOS}
			
				\begin{compactdesc}
					\item[BDAG 447a] \textbf{}
						\begin{compactitem}
							\item[1] \textbf{pertaining to being attractive in outward appearance, \textit{beautiful, handsome, fine}} in outward appearance
							\item[2] \textbf{pertaining to being in accordance at a high level with the purpose of something or someone, \textit{good, useful}}
							\item[2.a] of things
							\item[2.b] of moral quality (opposite \gr{αἰσχρός}) \textit{good, noble, praiseworthy, contributing to salvation}
								\begin{compactitem}
									\item \hyperref[JAMES-4-17]{James 4:17} is listed here.
									\item This use often appears with \gr{ἔργον}; BDAG lists many examples.
								\end{compactitem}
							\item[2.c] in any respect \textit{unobjectionable, blameless, excellent}
							\item[2.c.$\alpha$] of persons
							\item[2.c.$\beta$] of things
							\item[2.d] the term \gr{καλόν} (\gr{ἐστιν}) in the general sense\textit{it is good} qualifies items that fit under one of the preceding classifications
							\item[2.d.$\alpha$] \textit{it is pleasant, desirable, advantageous}
							\item[2.d.$\beta$] \textit{it is morally good, pleasing to God, contributing to salvation}
							\item[2.d.$\gamma$] (constructions with \gr{μᾶλλον})
						\end{compactitem}
					\item[CGL 739b, PDF 779] \textbf{}
						\begin{compactitem}
							\item[1] \textbf{beautiful, handsome, good-looking}
							\item[2] (of places, features of the natural world) \textbf{beautiful, fair}
							\item[3] (of created things, such as clothes, weapons, buildings) \textbf{beautiful, handsome, fine}
							\item[4] (of a voice, singing or similar) \textbf{beautiful}
							\item[5] (as a term of general commendation, of things or circumstances) good (in terms of quality, practice usefulness or capacity to satisfy or give pleasure), \textbf{good, excellent, fine}
							\item[6] (specifically, of a wind) \textbf{favourable}; (of sacrifices or omens) \textbf{good, auspicious, propitious}; (of coinage) \textbf{genuine}
							\item[7] neuter impersonal, it is a good time ---with infinitive \textit{to do something}
							\item[8] (prepositional phrase) \textit{in a good place, in the right spot}
							\item[9] (with moral connotation, of things said or done) \textbf{good, noble, honourable, fitting, proper}
							\item[10] (phrase, of a man combining moral excellence with good birth or social status, frequently as substantive) \textit{fine and good, upstanding, decent}
						\end{compactitem}
					\item[LSJ 870a, PDF 941] \textbf{}
						\begin{compactitem}
							\item[A.I.1] \textit{beautiful} of outward form
							\item[A.I.2.a] in Attic added to a name in token of love or admiration
							\item[A.I.2.b] as epithet
							\item[A.I.2.c] of divinities
							\item[A.I.3] \textit{beauty} (neuter substantive)
							\item[A.II.1] with reference to use, \textit{good, of fine quality}
							\item[A.II.2] of sacrifices, \textit{auspicious}
							\item[A.III.1] in a moral sense, \textit{beautiful, noble, honourable}
							\item[A.III.2] (neuter substantive) \textit{moral beauty, virtue, honour}, opposite \gr{αἰσχρόν}
							\item[A.III.3] of persons, in early writers coupled with \gr{ἀγαθός}
							\item[A.IV] in Attic and Tragedy frequently ironically, \textit{fine, specious}
							\item[B] degrees of comparison
							\item[C.I] adverbial:---Poets frequently use neuter \gr{καλόν} as adverb
							\item[C.II] regular adverb \gr{καλῶς}
								\begin{compactitem}
									\item Many definitions for the adverb \gr{καλῶς} follow here.
								\end{compactitem}
							\item[D] (certain compounds)
							\item[E] (discussion of syllabic quantity for poetry)
						\end{compactitem}
				\end{compactdesc}
				
				% \begin{compactdesc}
					% \item[Some NT uses] \textbf{}
						% \begin{compactdesc}
							% \item[]
						% \end{compactdesc}
				% \end{compactdesc}
				
			%
			\subparagraph{TDNT: \texorpdfstring{\gr{καλός}}{} (3:536--556)}
			%
				\label{KALOS-TDNT}
				
				``\gr{κινδυνεύει κατὰ τὴν ἀρχαίαν παροιμίαν τὸ καλὸν φίλον εἶναι}, Plato \textit{Lysis} 216c.%
					\footnote{%
						In my complete works translation: ``Maybe the old proverb is right, and the beautiful is a friend'' (PDF 700).
						}
				This ancient Greek proverb adopted by Plato brings out the significance of the term \gr{καλός} for human life'' (536).
		
	% -----
	\subsubsection{Book of Mormon} % *BOM
	% -----
		\label{GOODNESS-BOM}
	
		% \begin{tabular}{ll}
			% Total occurrences in the BOM & 
		% \end{tabular}
		
		\begin{compactdesc}
			\item[{\hyperref[3NEPHI-4-5]{3 Nephi 4:5}}] \phantom{.}
				\begin{compactitem}
					\item This is the nature of evil: to despoil and ruin, and not to grow.
					\item Goes with verse 14 too, the end of evil.
					\item And verses 19--20, 29, etc.
					\item Compare what happens to the robbers who want to repent in \hyperref[3NEPHI-6-3]{3 Nephi 6:3}.
					\item Compare \hyperref[MORMON-4-4]{Mormon 4:4--5}
				\end{compactitem}
		\end{compactdesc}
		
	% -----
	\subsubsection{Doctrine and Covenants} % *DC
	% -----
		\label{GOODNESS-DC}
	
		% \begin{tabular}{ll}
			% Total occurrences in the DC & 
		% \end{tabular}
		
		\begin{compactdesc}
			\item[{\hyperref[DC98-3]{DC 98:3}}]
		\end{compactdesc}
		
	% % -----
	% \subsubsection{Pearl of Great Price} % *PGP
	% % -----
		% \label{GOODNESS-PGP}
	
		% \begin{tabular}{ll}
			% Total occurrences in the PGP & 
		% \end{tabular}
		
		% \begin{compactdesc}
			% \item[]
		% \end{compactdesc}
		
% % ----------------------
% \subsection{Key Phrases} % *KEYPHRASES *PHRASES
% % ----------------------
	% \label{GOODNESS-PHRASES}

	% \begin{compactdesc}
		% \item[] \textbf{\#times} | list
			% % \begin{compactdesc}
				% % \item[Bible anchor verses] \phantom{.}
					% % \begin{compactdesc}
						% % \item[Romans 5:5] \gr{}
					% % \end{compactdesc}
				% % \item[Commentary] ---
			% % \end{compactdesc}
	% \end{compactdesc}
	
% % ----------------------
% \subsection{Bible Dictionaries} % *DICTIONARIES
% % ----------------------
	% \label{GOODNESS-BD}

% ----------------------
\subsection{Restoration Usage} % *RESTORATION
% ----------------------
	\label{GOODNESS-RESTORATION}

	% % -----
	% \subsubsection{Joseph Smith} % *JS
	% % -----
		% \label{GOODNESS-JS}
	
		% \begin{compactdesc}
			% \item[{\hyperref[KEY]{Date/Source}}]
		% \end{compactdesc}
	
	% -----
	\subsubsection{Brigham Young} % *BY
	% -----
		\label{GOODNESS-BY}
	
		\begin{compactdesc}
			\item[{\hyperref[EMS-1860-BY-02-FEB-1860]{2 February 1860}}] In the evening the President discoursed to a few present upon the subject, that goodness resulted in eternal lives; and wickedness gradually became weaker and smaller until it was lost in element.
		\end{compactdesc}
	
% ----------------------
\subsection{Commentary and Studies} % *COMMENTARY *STUDIES
% ----------------------
	\label{GOODNESS-COMMENTARY}

	% % -----
	% \subsubsection{Key Points}
	% % -----
		% \label{GOODNESS-KEYS}
	
		% \begin{compactitem}
			% \item
		% \end{compactitem}
		
	% -----
	\subsubsection{The pragmatic gospel: goodness and truth}
	% -----
		\label{GOODNESS-truth}
		
		Remember that you can't separate this stuff from the epistemology stuff, it's all bound together. But the following is a list of verses that have stuck out to me particularly regarding goodness stuff---see special note of these verses (for all of them, check the actual verses for other notes and connections):
		
		\begin{compactdesc}
		
			\item[{\hyperref[MATTHEW-7-11]{Matthew 7:11}}] The Father knows how to give \textit{good} gifts and \textit{good} things to those who are asking.
			\item[{\hyperref[MATTHEW-7-15]{Matthew 7:15--20}}] Here, you're trying to judge who the \gr{ψευδοπροφήτης} is, and you do it by looking at whether a person is a \gr{ἀγαθός} tree who brings forth \gr{καλός} fruit.
			\item[{\hyperref[MATTHEW-12-10]{Matthew 12:10--12}}] Especially verse 12, ``doing well.'' Compare to the verse about sin in James?
			\item[{\hyperref[MATTHEW-12-33]{Matthew 12:33--35}}] 
			\item[{\hyperref[MATTHEW-13-1]{Matthew 13:1--9}}] The parable of the sower---note the quality of the ground especially in verse 8 (and note Christ's explanation of this parable later in the chapter, where he brings this up again---in fact, there's a huge amount of ``goodness'' talk throughout this chapter; see also the use of a form of \gr{σαπρός} is verse 48). See also \hyperref[MARK-4-3]{Mark 4:3--9} and \hyperref[LUKE-8-4]{Luke 8:4--8}. Think about Alma 32 in light of these verses.
			\item[{\hyperref[LUKE-3-8]{Luke 3:8--9}}]
			\item[{\hyperref[LUKE-6-6]{Luke 6:6--11}}]
			\item[{\hyperref[LUKE-6-43]{Luke 6:43--45}}]
			\item[{\hyperref[LUKE-7-31]{Luke 7:31--35}}]
			\item[{\hyperref[JOHN-3-17]{John 3:17--21}}] Check out the Greek words for ``evil,'' and note that it talks about ``doing the truth'' in verse 21. Some interesting stuff here.
			\item[{\hyperref[ROMANS-2-4]{Romans 2:4}}] \phantom{.}
				\begin{compactitem}
					\item This isn't really about truth/goodness, but it does seem a somewhat pragmatic attitude for God to have.
				\end{compactitem}
			\item[{\hyperref[ROMANS-2-25]{Romans 2:25--29}}] \phantom{.}
				\begin{compactitem}
					\item These verses are more closely related to the truth as goodness thing, though there's still some basic pragmatist stuff here too.
				\end{compactitem} 
			\item[{\hyperref[ROMANS-3-7]{Romans 3:7--8}}]
			
			\item[{\hyperref[1NEPHI-6-6]{1 Nephi 6:6}}] \phantom{.}
				\begin{compactitem}
					\item The scriptures contain things of ``worth.''
				\end{compactitem}
			\item[{\hyperref[1NEPHI-19-23]{1 Nephi 19:23}}] ``profit and learning'' (note the other times it occurs)
			\item[{\hyperref[2NEPHI-2-5]{2 Nephi 2:5}}] People are ``instructed sufficiently'' to make a basic judgment: discerning between \textit{good and evil}, not between something like ``true'' and ``false.'' Why would this be?
				\begin{compactitem}
					\item It could be that the choice between good and evil is the choice that really matters most.
					\item It could be that the choice between good and evil is much easier to make than the one between true and false (a lot of the time).
				\end{compactitem}
			\item[{\hyperref[2NEPHI-9-28]{2 Nephi 9:28}}]
			\item A few verses in 2 Nephi 26: 
				\begin{compactitem}
					\item \hyperref[2NEPHI-26-24]{Verse 24}
					\item \hyperref[2NEPHI-26-28]{Verse 28}
					\item \hyperref[2NEPHI-26-33]{Verse 33}
				\end{compactitem}
			\item[{\hyperref[2NEPHI-33-4]{2 Nephi 33:4}}]
			\item[{\hyperref[JACOB-5]{Jacob 5}}] \phantom{.}
				\begin{compactitem}
					\item They judge all the fruit in terms of whether it's good or bad.
				\end{compactitem}
			\item \hyperref[OMNI-1-25]{Omni 1:25}
			\item \hyperref[MOSIAH-4-11]{Mosiah 4:11--12}---here I'm thinking of trying to tie to Alma 32, since at the beginning of this verse it talks about coming to a knowledge of the glory of God, but then seems to explain what that means in terms of having had certain experiences.
				\begin{compactitem}
					\item Note also the connection here with Moroni 10:6 and ``just and true,'' which is another important thing to bring up about this verse.
				\end{compactitem}
			\item[{\hyperref[ALMA-18-34]{Alma 18:34}}] \phantom{.}
				\begin{compactitem}
					\item Mentions what is ``just and true.''
				\end{compactitem}
			\item[{\hyperref[ALMA-32-27]{Alma 32:27--43}}] \phantom{.}
				\begin{compactitem}
					\item Note that this sermon is an expansion of the parable of the sower, so make sure to compare with that parable and Christ's explanation of it to see what Alma is doing. \hyperref[LUKE-8-4]{Luke 8:4--18}, 
				\end{compactitem}
			\item[{\hyperref[HELAMAN-7-26]{Helaman 7:26}}]
			\item[{\hyperref[HELAMAN-12-26]{Helaman 12:26}}] \phantom{.}
				\begin{compactitem}
					\item It's about \textit{doing}, not believing.
					\item Also verse 24, before this one.
				\end{compactitem}
			\item \hyperref[ETHER-3-4]{Ether 3:4}, the ``benefit'' stuff
			\item \hyperref[ETHER-4-11]{Ether 4:11--12}, golden stuff
			\item[{\hyperref[ETHER-8-26]{Ether 8:26}}]
			\item[{\hyperref[MORONI-7-15]{Moroni 7:5--26}}] Verses 5--19 are the core of this passage, but the discussion goes on for a little longer. 
			\item \hyperref[MORONI-10-6]{Moroni 10:6}, obviously
			\item \hyperref[MORONI-10-18]{Moroni 10:18}; see also 10:25 here; also 10:30	
			
			\item[{\hyperref[DC6-33]{DC 6:33}}] 
			\item[{\hyperref[DC10-20]{DC 10:20}}] 
			\item[{\hyperref[DC11-12]{DC 11:12}}]
			\item[{\hyperref[DC18-18]{DC 18:18}}] Very important verse about the Spirit manifesting to us what is ``expedient.''
			\item[{\hyperref[DC46]{DC 46}}], the discussion of spiritual gifts, they are given for ``profit'' and ``benefit,'' not necessarily truth; worth looking into at some point, or using in favor of some argument
			\item[{\hyperref[DC50-22]{DC 50:22--23}}] Compare to Alma 32, especially the verse which says something like, ``it is light, because it is discernible,'' or whatever---the ``darkness'' in verse 23 here in DC 50 is a foil to that ``light.''
			\item[{\hyperref[DC56-2]{DC 56:2--4}}] Especially verse 4, ``as it seemeth me good.'' 
				\begin{compactitem}
					\item In this whole section really, notice how the Lord is revoking commandments in response to different things. Why didn't he just tell them what to do using his all-knowing powers? Why didn't he just tell them the final thing that they would end up doing? This is relevant for many sections in this middle part of the DC.
				\end{compactitem}
			\item[{\hyperref[DC58-25]{DC 58:25}}] (This one involves people ``counseling'' between themselves and the Lord to figure out what is best to do.)
			\item[{\hyperref[DC58-26]{DC 58:26--33}}] \phantom{.}
				\begin{compactitem}
					\item We just need to be out there, doing good.
				\end{compactitem} 
			\item[{\hyperref[DC67-9]{DC 67:9}}] relationship between righteousness and goodness?
			\item[{\hyperref[DC68-4]{DC 68:4}}] pragmatic stuff?
			\item[{\hyperref[DC81-4]{DC 81:4}}] pragmatic stuff, including talk of the ``greatest good''
			\item[{\hyperref[DC82]{DC 82}}] there is a lot about goodness and things turning out in the right way in this section; 8--9 is good, but I think there are others. 16 for example, or 15--16.
			\item[{\hyperref[DC84-73]{DC 84:73}}]
			\item JD 14:152, from BY, talking about the name ``Mormon''; there is a way of seeing the ``more good'' thing not as the meaning of the name Mormon (although JS thought of it that way), and seeing it more as the slogan captured in the name ``Mormon'' or ``Mormonism.'' Very interesting idea.
			\item[{\hyperref[DC97]{DC 97:7--9}}]
			\item[{\hyperref[DC98-11]{DC 98:11}}] 
			
			\item Note that, in the meeting where they discussed publishing the Book of Commandments, those who signed the ``Testimony'' about them described them as ``profitable for all men.'' See that document \hyperref[EMS-1831-JS-CIR-2-NOV-1831-TEST]{here}.
			
		\end{compactdesc}